% =============================================================================
%  bookstyle.tex: the one visual vocabulary shared by every figure in the book
% =============================================================================
%  Loaded by preamble.tex (and by preview.tex for drawing a figure on its own).
%  Every TikZ/PGFPlots figure in figures-src/tikz/ draws with these names only:
%  no raw HTML colours, no ad hoc line widths, no font sizes set by hand.
%
%  Rules of thumb
%   * Draw at the printed size: never scale a figure with width= or scale=
%     after the fact, so text stays at the sizes set here.  Labels are \small
%     (9.6pt with the 10.5pt body), which puts their subscripts at 6.7pt, above the
%     book's 6.5pt floor;
%     \footnotesize is the floor for plain text and must not carry subscripts.
%   * Colour encodes meaning, and the meaning is the same in every chapter
%     (the book-wide legend; a figure may depart from it only to pair or
%     distinguish several vectors of one kind, e.g. third-law pairs or Lami's
%     three forces, and then the caption says so):
%       crimson  every force: weight, normal reaction, tension, gravity, F
%       accent   friction and applied pushes/pulls, where they must be told
%                apart from the contact forces; kinetic-energy bars (K)
%       forest   velocity v and angular velocity omega
%       primary  acceleration a and angular acceleration alpha; main plot curve
%       teal     momentum p and angular momentum L; components and legs
%       plum     potential energy U, total energy E, potentials
%       steel    dimensions, construction lines, secondary labels
%       grey     guides, paths not taken, ghost positions
%   * Labels sit beside what they name, in the colour of what they name.
% -----------------------------------------------------------------------------

% ---------------------------------------------------------------- palette ---
% Prefixed with c so they never collide with the text greys in preamble.tex
% (where `ink' is pure black for the prose).
\definecolor{cInk}{HTML}{102A43}     % outlines, axes, neutral labels
\definecolor{cPrim}{HTML}{1F6FEB}    % blue
\definecolor{cTeal}{HTML}{0E7C7B}
\definecolor{cAcc}{HTML}{E85D04}     % orange
\definecolor{cCrim}{HTML}{C1121F}    % red
\definecolor{cFor}{HTML}{2D6A4F}     % green
\definecolor{cSteel}{HTML}{486581}
\definecolor{cPlum}{HTML}{7B2CBF}
\definecolor{cAmber}{HTML}{F4A261}
\definecolor{cGrey}{HTML}{9FB3C8}
% fills (pale, so outlines and arrows read on top of them)
\definecolor{cEarth}{HTML}{BFD3E6}   % planets
\definecolor{cMoon}{HTML}{D9E2EC}    % small bodies, light blocks
\definecolor{cBlock}{HTML}{93C5C6}   % blocks, bodies under study
\definecolor{cSlope}{HTML}{E1E7EF}   % inclines, ground fill
\definecolor{cRock}{HTML}{9AA8B6}    % mountains, walls, fixed supports
\definecolor{cSun}{HTML}{F4A261}

% ------------------------------------------------------------------ lines ---
\tikzset{
  % line weights, used everywhere instead of numbers
  outline/.style={line width=0.8pt,draw=cInk},
  thin line/.style={line width=0.5pt},
  hair line/.style={line width=0.35pt},
  curve/.style={line width=1.3pt},
  % ------------------------------------------------------------- arrows -----
  % one arrow head for the whole book; size scales with the line it ends
  bookhead/.tip={Stealth[length=2.4mm,width=1.7mm]},
  smallhead/.tip={Stealth[length=1.8mm,width=1.3mm]},
  vector/.style={-bookhead,line width=1.1pt,line cap=round},
  force/.style={vector,cCrim},
  applied/.style={vector,cAcc},
  velocity/.style={vector,cFor},
  accel/.style={vector,cPrim},
  momentum/.style={vector,cTeal},
  % dimension lines and leaders
  dim/.style={{smallhead}-{smallhead},thin line,cSteel},
  dimtick/.style={thin line,cSteel,densely dashed},
  leader/.style={hair line},
  guide/.style={thin line,cGrey,dash pattern=on 3pt off 2pt},
  orbitpath/.style={thin line,cGrey,dash pattern=on 4pt off 3pt},
  % --------------------------------------------------------------- bodies ---
  body/.style={fill=cBlock,draw=cSteel,line width=0.6pt},
  planet/.style={fill=cEarth,draw=cInk,line width=0.9pt},
  moon/.style={fill=cMoon,draw=cInk,line width=0.7pt},
  point mass/.style={circle,fill=cInk,inner sep=0pt,minimum size=3.2pt},
  % `ground' hatches a CLOSED region (a thin rectangle under the surface);
  % on a bare line it draws only the line.  For the usual pale strip under a
  % surface use `ground strip' on a rectangle, then draw the surface with `outline'.
  ground/.style={line width=0.8pt,draw=cInk,
    postaction={draw=none,pattern={Lines[angle=45,distance=2.2pt,line width=0.35pt]},
      pattern color=cInk}},
  ground strip/.style={fill=cSlope,draw=none},
  % a thin curved arrow for omega, alpha and other senses of rotation
  rotation/.style={-smallhead,thin line,cInk},
  % a small square marking a right angle; use as  \draw[right angle] (corner) ++(...)
  right angle/.style={hair line,cInk},
  % a medium-weight construction line in any colour: [construction,cFor]
  construction/.style={line width=0.8pt},
  % a dot in any colour: \node[dot,cCrim] at (P) {};
  dot/.style={circle,fill,inner sep=0pt,minimum size=3.2pt},
  % angle arcs: \draw[angle arc] (A) arc[...];  label it with a plain node
  angle arc/.style={thin line,cInk},
  % Riemann strips and shaded areas: pale fill, darker hairline edge
  strip/.style={fill=cPrim!14,draw=cPrim!60,hair line},
  area fill/.style={fill=cPrim!12,draw=none},
  % components of a vector (dashed, same colour as the vector, lighter head)
  component/.style={-smallhead,thin line,dash pattern=on 3pt off 1.5pt},
  % a resultant built by parallelogram: dashed, full weight
  resultant/.style={dash pattern=on 4pt off 2pt},
  parallelogram/.style={hair line,cGrey,dash pattern=on 2pt off 1.5pt},
  % rotation axes and hidden (back) edges of solids
  rotation axis/.style={thin line,cInk,dash pattern=on 6pt off 2pt on 1pt off 2pt},
  hidden/.style={dash pattern=on 2pt off 1.5pt},
  % the slice in a derivation figure (strip, ring, disc, shell)
  slice/.style={fill=cAcc!40,draw=cAcc,thin line},
  % a ghost: an earlier or later position of the same body.  Opacity also fades
  % any node placed on the same path, so put a ghost's label in a separate node.
  ghost/.style={opacity=0.35},
  % a particle or bob: shaded ball with a steel rim;  \pic at (P) {particle};
  particle/.pic={\shade[ball color=cBlock] (0,0) circle[radius=0.11];
                 \draw[cSteel,thin line] (0,0) circle[radius=0.11];},
  % stable and unstable equilibrium points on a U(x) curve
  stable point/.style={circle,fill=cInk,inner sep=0pt,minimum size=4pt},
  unstable point/.style={circle,draw=cInk,fill=white,thin line,inner sep=0pt,minimum size=4pt},
  % energy bars
  U bar/.style={fill=cPlum!55,draw=cPlum,thin line},
  K bar/.style={fill=cAcc!55,draw=cAcc,thin line},
  % horizontal circles drawn in perspective use this y-squash for the ellipse
  % (x radius r, y radius \persp*r)
  % the centre-of-mass marker: \pic at (G) {cm};
  cm/.pic={\fill[white] (0,0) circle[radius=0.09];
           \fill[cInk] (0,0) -- (0.09,0) arc[start angle=0,end angle=90,radius=0.09] -- cycle;
           \fill[cInk] (0,0) -- (-0.09,0) arc[start angle=180,end angle=270,radius=0.09] -- cycle;
           \draw[cInk,hair line] (0,0) circle[radius=0.09];},
  % out of / into the page
  out of page/.pic={\draw[cInk,thin line,fill=white] (0,0) circle[radius=0.09];
                    \fill[cInk] (0,0) circle[radius=0.025];},
  into page/.pic={\draw[cInk,thin line,fill=white] (0,0) circle[radius=0.09];
                  \draw[cInk,thin line] (-0.064,-0.064) -- (0.064,0.064) (-0.064,0.064) -- (0.064,-0.064);},
  % caution: `knockout' fills with the CURRENT colour if a bare colour name is
  % given alongside it; colour knocked-out text with text=<colour>.
  % the contact point of a rolling body
  contact point/.style={circle,fill=cInk,inner sep=0pt,minimum size=2.6pt},
  % ---------------------------------------------------------------- text ----
  % every figure starts from \small (9.6pt; subscripts 6.7pt); \footnotesize (8.75pt)
  % is the floor, for plain words only
  figure text/.style={font=\small,text=cInk},
  small label/.style={font=\footnotesize},
  note/.style={font=\small\itshape,text=cInk},
  % a label that must sit over a line: white knock-out behind the text
  knockout/.style={fill=white,inner sep=1.2pt},
  % applied to every figure by \tikzfig
  % default drawing colour is ink, so plain `dot', `fill' and `draw' match the palette
  bookfigure/.style={figure text,color=cInk,line join=round,>=bookhead},
}

\def\persp{0.25}
% ----------------------------------------------------------------- plots ----
\usepgfplotslibrary{groupplots}
\pgfplotsset{
  % the default look of every plot: two axis lines with arrow tips, no box,
  % no grid; ticks only where the text refers to a value
  bookplot/.style={
    axis lines=left,
    axis line style={line width=0.7pt,draw=cInk,-{Stealth[length=2mm,width=1.5mm]}},
    every tick/.style={line width=0.5pt,draw=cInk},
    tick label style={font=\small,text=cInk},
    label style={font=\small,text=cInk},
    every axis plot/.append style={curve},
    xtick=\empty, ytick=\empty,
    clip=false,
    samples=120,
    enlargelimits=false,
    axis on top=false,
  },
  % a zero line drawn inside a plot (energy plots)
  zero line/.style={thin line,cInk},
}
